% Options for packages loaded elsewhere
\PassOptionsToPackage{unicode}{hyperref}
\PassOptionsToPackage{hyphens}{url}
\documentclass[
]{article}
\usepackage{xcolor}
\usepackage{amsmath,amssymb}
\setcounter{secnumdepth}{-\maxdimen} % remove section numbering
\usepackage{iftex}
\ifPDFTeX
  \usepackage[T1]{fontenc}
  \usepackage[utf8]{inputenc}
  \usepackage{textcomp} % provide euro and other symbols
\else % if luatex or xetex
  \usepackage{unicode-math} % this also loads fontspec
  \defaultfontfeatures{Scale=MatchLowercase}
  \defaultfontfeatures[\rmfamily]{Ligatures=TeX,Scale=1}
\fi
\usepackage{lmodern}
\ifPDFTeX\else
  % xetex/luatex font selection
\fi
% Use upquote if available, for straight quotes in verbatim environments
\IfFileExists{upquote.sty}{\usepackage{upquote}}{}
\IfFileExists{microtype.sty}{% use microtype if available
  \usepackage[]{microtype}
  \UseMicrotypeSet[protrusion]{basicmath} % disable protrusion for tt fonts
}{}
\makeatletter
\@ifundefined{KOMAClassName}{% if non-KOMA class
  \IfFileExists{parskip.sty}{%
    \usepackage{parskip}
  }{% else
    \setlength{\parindent}{0pt}
    \setlength{\parskip}{6pt plus 2pt minus 1pt}}
}{% if KOMA class
  \KOMAoptions{parskip=half}}
\makeatother
\usepackage{longtable,booktabs,array}
\usepackage{calc} % for calculating minipage widths
% Correct order of tables after \paragraph or \subparagraph
\usepackage{etoolbox}
\makeatletter
\patchcmd\longtable{\par}{\if@noskipsec\mbox{}\fi\par}{}{}
\makeatother
% Allow footnotes in longtable head/foot
\IfFileExists{footnotehyper.sty}{\usepackage{footnotehyper}}{\usepackage{footnote}}
\makesavenoteenv{longtable}
\setlength{\emergencystretch}{3em} % prevent overfull lines
\providecommand{\tightlist}{%
  \setlength{\itemsep}{0pt}\setlength{\parskip}{0pt}}
\usepackage{bookmark}
\IfFileExists{xurl.sty}{\usepackage{xurl}}{} % add URL line breaks if available
\urlstyle{same}
\hypersetup{
  hidelinks,
  pdfcreator={LaTeX via pandoc}}

\author{}
\date{}

\begin{document}

\section{}\label{section}

\section{SÍLABO DOSIFICADO DE MÉTODOS NUMÉRICOS PARA ECUACIONES
DIFERENCIALES PARCIALES
(CM032)}\label{suxedlabo-dosificado-de-muxe9todos-numuxe9ricos-para-ecuaciones-diferenciales-parciales-cm032}

Profesor: Fidel Jara Huanca.

\begin{longtable}[]{@{}
  >{\raggedright\arraybackslash}p{(\linewidth - 8\tabcolsep) * \real{0.2000}}
  >{\raggedright\arraybackslash}p{(\linewidth - 8\tabcolsep) * \real{0.2000}}
  >{\raggedright\arraybackslash}p{(\linewidth - 8\tabcolsep) * \real{0.2000}}
  >{\raggedright\arraybackslash}p{(\linewidth - 8\tabcolsep) * \real{0.2000}}
  >{\raggedright\arraybackslash}p{(\linewidth - 8\tabcolsep) * \real{0.2000}}@{}}
\toprule\noalign{}
& & & & \\
\midrule\noalign{}
\endhead
\bottomrule\noalign{}
\endlastfoot
SEM & \textbf{SESESIÓN 1 (2h Teoría)} & \textbf{SESESIÓN 2 (2h
Laboratorio)} & SESIÓN 3 (2h Teoría) & \\
1\textbf{} & Entrega del sílabo. Introducción al curso. Prueba de
entrada. & Introducción al software Python. Esquema de diferencias
finitas para un problema de valor de frontera. & Aproximación vía
diferencias finitas hasta solución generalizada. {[}1{]} & \\
2 & Consistencia, convergencia y estabilidad. Teorema de equivalencia de
Lax. Orden de convergencia. {[}1{]} & Esquemas explícitos e implícitos
para la ecuación de calor y ejercicios sugeridos. {[}1{]} & Operadores
de dos niveles. Consistencia y estimación de error. Ejemplos. {[}1{]}
& \\
3 & Métodos de diferencias finitas para problemas elípticos. El problema
de Poisson en una dimensión. Consistencia, estabilidad y convergencia.
{[}2{]} & Practica Calificada 1. & Problema de Poisson en dos
dimensiones. Consistencia, estabilidad y convergencia. {[}2{]} & \\
4 & Método de Galerkín. Lema de Cea. Teorema de Lax Milgram. Base de
Lagrange. {[}2{]} & Esquema de 5 puntos para la ecuación de Poisson
unidimensional y Esquema de 9 puntos para la ecuación de Poisson
bidimensional. Solución por métodos iterativos. {[}3{]} & Construcción
de la matriz de rigidez y del vector de carga en el método de Galerkin.
{[}2{]} & \\
5\textbf{} & Motivación. Ecuación de transporte. Ondas acústicas. Flujo
de tráfico. Aguas poco profundas. Ecuaciones de Euler. {[}4{]} &
Practica Calificada 2 & Leyes de conservación del tipo hiperbólicos
escalar y sistema de leyes de conservación. Solución clásica, débil y
entrópica. Ondas de choque. {[}4{]} & \\
6 & Condición Courant-Friedrichs-Lewy. Ecuación de Transporte. Ecuación
de Burgers. Consistencia. Estabilidad de Von-Neumann

Esquemas Upwind, FTCS, Lax-Friedrichs, Leapfrog. El problema de Riemman.
{[}5{]} & Ejercicios sugeridos:

Ejercicios 3.1, 3.2, 3.3, 3.4, 3.5, 3.6, 3.7, 3.8 {[}6{]} & Esquema
BTCS, Lax-Wendroff, Crank-Nicolson, Du Fort-Frankel. {[}5{]} & \\
7 & Ecuaciones de Agua Poco Profundas unidimensional. Problema de Dam
Break.

Esquema de Mac-Cormack. {[}3{]} & Ejercicios sugeridos:

Detallar ejemplos de la sección 5.1 {[}7{]}

Ejercicios del capitulo 29 {[}2{]} & Presentación del Proyecto. Parte 1
& \\
8 & & EXAMEN PARCIAL & & \\
9 & Propiedad monótona de esquemas: Teorema de Godunov {[}6{]} &
Tutorial de Clawpack & Método de Volúmenes Finitos.

Problema de Dirichlet. Estimación de error. & \\
10\textbf{} & Problemas parabólicos e hiperbólicos. {[}4{]} & Practica
calificada 3 & Esquema numérico para el caso lineal y no lineal. {[}4{]}
& \\
11\textbf{} & Método de Gudonov y variantes {[}4{]} & Ejercicios
sugeridos & El teorema de Lax Wendroff {[}6{]} & \\
12 & Resolvente de Riemann {[}4{]} aproximado & Practica calificada 4 &
Método de direcciones {[}6{]} alternadas & \\
13 & Método de Weno y Eno {[}6{]} & Ejercicios sugeridos & Técnicas de
descomposición de flujo & \\
14\textbf{} & Esquema de variación total decreciente y consistente con
entropía {[}6{]} & Ejercicios sugeridos & Métodos de volúmenes finitas
en mallas no estructuradas & \\
15\textbf{} & Aplicaciones {[}4{]} & Practica calificada 5 &
Presentación del Proyecto. Parte II & \\
16 & & & & \\
17 & & & & \\
\end{longtable}

{[}1{]} Theoretical Numerical Analysis: A Functional Analysis Framework.
Kendall Atkinson.

{[}2{]} Classical Numerical Analysis: A Comprehensive Course. Abner J.
Salgado.

{[}3{]} Numerical Methods for Partial Differential Equations. Jon Shiach
\href{https://jonshiach.github.io/teaching/PDEs}{\emph{\emph{https://jonshiach.github.io/teaching/PDEs}}}

{[}4{]} Solving Hyperbolic Equations with Finite Volume Methods. María
Elena Vázquez-Cendón.

{[}5{]} El método de las diferencias finitas. Luis Lara Romero, Zenner
Chávez Aliaga, José Castañeda Vergara.
\href{https://static.upao.info/descargas/7785b78834b257ebc32c5fdbbbeb51844110d47e0c1b7621073297f7054588f8feabedcf65f22243e1c94da5fd1fdf63312e1c80f31ef101692d7a1e07e3ba49/el-mEtodo-de-diferencias-finitas.pdf}{\emph{\emph{https://static.upao.info/descargas/7785b78834b257ebc32c5fdbbbeb51844110d47e0c1b7621073297f7054588f8feabedcf65f22243e1c94da5fd1fdf63312e1c80f31ef101692d7a1e07e3ba49/el-mEtodo-de-diferencias-finitas.pdf}}}

{[}6{]} Numerical Methods for Conservation Laws. Randall J. LeVeque

{[}7{]} Numerical Methods for Conservation Laws. Jan S. Hesthaven

\end{document}
